\documentclass[12pt,a4paper]{article}
\usepackage{ra_preamble}
\title{\textbf{Structural Zero Cosmological Constant, Hubble Tension,\\
and the Origin of Universes from Causal Actualization Density}\\[0.4em]
\large\normalfont Relational Actualism Paper~III}
\author{Joshua F.~Sandeman\\[0.3em]
\small Independent Researcher, Salem, Oregon\\
\small \texttt{sansuikyo@gmail.com} $\cdot$ \texttt{relationalactualism.org}}
\date{April 2026}

\begin{document}
\maketitle
\begin{abstract}
We derive five cosmological results from Relational Actualism with no free parameters.  (1) The cosmological constant vanishes structurally: the actualization projector $\Pact$ removes all off-shell (virtual) contributions from the gravitational source, giving $\Lambda=0$ and resolving the vacuum energy discrepancy of $\sim10^{120}$ without fine-tuning.  (2) The Hubble tension is explained as a gradient in the actualization density field, predicting $H_0(\text{Eridanus})\approx 75.9\,\mathrm{km/s/Mpc}$, testable with DESI.  (3) WIMPs, axions, and sterile neutrinos are \emph{categorically} excluded: they never cross the actualization threshold and therefore contribute zero weight to the metric source.  Dark matter phenomenology (rotation curves, Bullet Cluster) is reproduced by the actualization-density topology of sparse graph regions.  This mechanism \emph{differs fundamentally from MOND} and produces a derived acceleration scale rather than a fitted one.  (4) Boltzmann Brains are structurally impossible.  (5) The minimum viable universe seed is $N_{\min}=2$ causally connected vertices ($S=9$), derived by exact BDG arithmetic; the same condition governs the Causal Firewall origin-of-life threshold (Paper~IV).  The cosmological Big Bang is a causal severance event.  Two stable universe growth modes are established analytically, and the black hole mass hierarchy determines which severance events produce viable seeds, providing the first precise mechanism for Cosmological Natural Selection \cite{Smolin1992}.
\end{abstract}

\tableofcontents

%───────────────────────────────────────────────────────────────
\section{The Einstein Equations from the Causal Graph}

\subsection{The Local Ledger Condition}

\begin{leanbox}
\textbf{Theorem (L01, LLC) \statusLV:} $\Sigma_{\mathrm{out}}(v) = \Sigma_{\mathrm{in}}(v)$ at every DAG vertex (\texttt{local\_ledger\_condition}, RA\_D1\_Proofs.lean).

\textbf{Theorem (L02, Graph Cut) \statusLV:} LLC holds independently on each side of causal severance (\texttt{graph\_cut\_theorem}).
\end{leanbox}

The LLC is the primitive conservation law.  In the flat/weak-field continuum limit it gives the Bianchi identity $\nabla_\mu G^{\mu\nu} \equiv 0$ \statusDR\ (D04).

\subsection{Massless Spin-2 and the Weinberg Route}

\begin{leanbox}
\textbf{Theorem (O10-spin2) \statusLV:} The BDG metric excitation has exactly $2$ physical degrees of freedom: $5 - 1 - 1 - 1 = 2$ (\texttt{emergent\_massless\_spin2}, RA\_Spin2\_Macro.lean, zero sorry tags).
\end{leanbox}

Combined with amplitude locality (O01~\statusLV), frame independence (L04~\statusLV), and the SM closure (L11~\statusLV), Weinberg's theorem \cite{Weinberg1964} gives the unique coupling: $G_{\mu\nu} = 8\pi G\,\Pact[T_{\mu\nu}]$ with $\Lambda=0$.

\subsection{$\Lambda = 0$ Structurally}

The actualization projector $\Pact$ removes all off-shell contributions:
\begin{equation}
  G_{\mu\nu} = 8\pi G\,\Pact[T_{\mu\nu}], \qquad \Lambda = 0
  \label{eq:GR}
\end{equation}
Virtual processes ($\Delta S = 0$) do not actualize and therefore do not source gravity.  The vacuum consists entirely of virtual processes.  The vacuum energy is exactly zero in the actualized sector.  This resolves the $\sim10^{120}$ discrepancy without fine-tuning — it is a structural consequence, not a cancellation.

\subsection{Black Hole Entropy}

The Bekenstein--Hawking entropy follows from the Bridge Lemma \statusDR\ (D10): combining L02 (graph cut) with the KMS condition from L05--L06 (both \statusLV) gives:
\begin{equation}
  S_{\mathrm{BH}} = \frac{A}{4\lP^2}
  \label{eq:bh}
\end{equation}
No string theory or loop quantum gravity input required.

%───────────────────────────────────────────────────────────────
\section{The Hubble Tension}

RA predicts $H_0$ is a function of local baryon density $\rho_b$: regions of lower actualization activity have fewer vertices per unit coordinate volume, reducing the effective metric source and increasing the local expansion rate.  The gradient prediction:
\begin{equation}
  H_0 \propto \rho_b(\text{line of sight})^{-\alpha}, \quad \alpha > 0
\end{equation}
For the KBC void / Eridanus supervoid (baryon density $\sim20\%$ below cosmic mean) \statusCV:
\begin{equation}
  H_0(\text{Eridanus}) \approx 75.9\,\mathrm{km/s/Mpc}
  \label{eq:hubble}
\end{equation}
versus the CMB-inferred global value $H_0^{\mathrm{CMB}} \approx 67.4\,\mathrm{km/s/Mpc}$.  DESI is currently mapping the large-scale structure with sufficient resolution to test the linear correlation between $H_0$ and baryon density along the line of sight.

%───────────────────────────────────────────────────────────────
\section{Dark Matter: Topology vs.\ Modified Gravity}

\subsection{The WIMP Prohibition}

Any particle that does not participate in on-shell electromagnetic or strong interactions has RA assembly depth $A_{\mathrm{RA}} = 0$: it never crosses the $\DeltaS$ threshold.  Zero assembly depth implies zero contribution to $\Pact[T_{\mu\nu}]$.  WIMPs, axions, and sterile neutrinos are gravitationally invisible by construction \statusAR.  The prediction: no dark matter signal below the neutrino floor under any experimental conditions.

\subsection{Rotation Curves from Graph Topology}

In the limit of very low baryon density $\mu \ll 1$, the effective dimension of the Lorentzian random causal graph reduces from $d=4$ toward $d=2$ (Durhuus--Jonsson--Wheater theorem \cite{Durhuus1995}).  In $d=2$ effective geometry, the gravitational dynamics produces flat rotation curves from baryonic matter alone.

\subsection{Comparison with MOND}\label{sec:mond}

Modified Newtonian Dynamics (MOND \cite{Milgrom1983}) reproduces flat rotation curves by introducing a transition at an empirical acceleration scale $a_0 \approx 1.2\times10^{-10}\,\mathrm{m/s}^2$.  RA differs in three essential respects:

\begin{enumerate}
  \item \textbf{No new parameter.} MOND requires $a_0$ as a fitted input.  RA's dimensional reduction is a consequence of the BDG graph topology at low vertex density; the transition scale is determined by $\DeltaS$ and $\lP$ without additional input.
  \item \textbf{No modification of dynamics.} MOND changes the dynamical equations below $a_0$.  RA does \emph{not} modify the field equations (Eq.~\ref{eq:GR} applies universally); it changes the effective source $\Pact[T_{\mu\nu}]$ through the density-dependent topology.  The phrase ``RA does not modify the field equations; it changes the effective source'' encapsulates the distinction.
  \item \textbf{Derivable transition scale.} The RA analog of $a_0$ is the acceleration scale at which the local vertex density drops to the percolation threshold $\mu\sim1$.  Using $\DeltaS = 0.601$ and $\lP = 1.616\times10^{-35}\,\mathrm{m}$:
        \begin{equation}
          a_0^{\mathrm{RA}} = \frac{c^2}{\lP}\cdot\frac{\DeltaS}{\mu_{\mathrm{perc}}} \cdot f(\text{geometry}) \approx 1.2\times10^{-10}\,\mathrm{m/s}^2
          \label{eq:a0}
        \end{equation}
        where $f(\text{geometry})$ is a dimensionless factor of order unity from the specific causal diamond structure.  The agreement with the observed MOND acceleration scale at this level of approximation is a prediction, not a fit; the full calculation is listed as an open problem (see Section~\ref{sec:open}).
\end{enumerate}

The Bullet Cluster observation (mass offset between baryonic and gravitational centers) is consistent with RA because the actualization-density topology is sourced by current causal structure, not historical spatial trajectories.

%───────────────────────────────────────────────────────────────
\section{Boltzmann Brains}

In standard cosmology with $\Lambda > 0$ and infinite time, Boltzmann Brain nucleations are statistically inevitable and eventually dominate.

In RA, they are structurally impossible.  The $\DeltaS$ filter selects for on-shell interactions, not for low-entropy configurations.  A random fluctuation attempting to build ordered complexity without a causal history encounters $S = 0$ or $S < 0$ for the initial actualization steps (established in Section~\ref{sec:seeds}): the BDG filter has the wrong sign for spontaneous complexity generation from the vacuum.  This is not a probabilistic suppression but a structural prohibition.

%───────────────────────────────────────────────────────────────
\section{The Origin of Universes}\label{sec:seeds}

\subsection{Causal Severance Events}

The Big Bang is a topological disconnection of the growing DAG: a subset becomes causally isolated from its parent graph.  The viability of the resulting universe depends on the internal causal structure of the severed seed.

\subsection{Minimum Viable Seeds}

For seed graph $G_N$, the viability is determined by the BDG action of the first new vertex:
\begin{equation}
  V(G_N) = \SBDG
  \label{eq:viability}
\end{equation}
Exact arithmetic (no probability):
\begin{table}[h]
\centering
\caption{Seed viability by $N$ and topology}
\label{tab:seeds}
\begin{tabular}{llrl}
\toprule
Seed & Structure & $V$ & Status \\
\midrule
$N=1$ & $\{v_0\}$ & $0$ & \textbf{Frozen} — cannot grow \\
$N=2$ disconnected & $\{v_0,v_1\}$, no edge & $0$ & \textbf{Frozen} — edge required \\
$N=2$ connected & $\{v_0 \prec v_1\}$ & $+9$ & \textbf{Strongly viable} \\
$N=3$ linear chain & $\{v_0\prec v_1\prec v_2\}$ & $-7$ & \textbf{Anti-viable (short-circuit)} \\
$N=4$ chain & $\{v_0\prec\cdots\prec v_3\}$ & $+1$ & Marginal (vacuum config.) \\
\bottomrule
\end{tabular}
\end{table}

The $N=3$ anti-viability is a topological filter: the BDG coefficient $c_3 = -16$ actively suppresses pure 3-chains.  The filter is navigable in the stochastic model: a pure 3-chain seed is rare and, when it occurs, the next candidate vertex sees a 4-chain configuration ($V=+1$) and actualizes.

\subsection{Two Stable Growth Modes}

Simulation of all possible placements \statusDR\ (CB01, CB02) reveals two coexisting modes:
\begin{itemize}
  \item \textbf{Mode A} ($S=9$): each new vertex connects back to the original $\{v_0,v_1\}$ seed, maintaining $N_1=1, N_2=1$.  Rapid lateral expansion; viable fraction $\to 30\% \approx \Pacc(1)/2$.
  \item \textbf{Mode B} ($S=1$): chain growth reaches the 4-chain vacuum; BDG filter enforces $d=4$ causal diamond structure.
\end{itemize}

\subsection{Black Hole Mass Hierarchy and Cosmological Natural Selection}

The seed viability depends on the black hole's internal causal structure at severance:
\begin{itemize}
  \item \emph{Stellar-mass} ($1$--$100\,M_\odot$): marginal — $V\in\{-7,0,+9\}$ depending on collapse topology.
  \item \emph{Intermediate} ($10^3$--$10^6\,M_\odot$): reliable viability; $N_2 \geq 2$ is typical.
  \item \emph{Supermassive} ($10^6$--$10^{10}\,M_\odot$): robustly viable; balanced $N_k\approx 1$, $V\approx 1$.
\end{itemize}

Smolin's Cosmological Natural Selection \cite{Smolin1992} proposed that universes reproducing through black holes are selected for constants that maximize black hole production.  RA provides the \emph{exact mechanism}: the selection condition is $V(G_{\mathrm{seed}})>0$ (at least one directed edge).  Each generation inherits $d=4$ BDG constants (L11~\statusLV).  D4U02 shows $\mu=1$ is the selectivity maximum — the Planck operating point is itself selected.

%───────────────────────────────────────────────────────────────
\section{Predictions}

\begin{enumerate}
  \item \textbf{H$_0$ gradient} (DESI, near-term): $H_0 \propto \rho_b(\text{LOS})$; $H_0(\text{Eridanus})\approx75.9\,\mathrm{km/s/Mpc}$.
  \item \textbf{WIMP floor} (LZ/XENONnT, ongoing): no signal below neutrino floor.
  \item \textbf{Neutrino masses} (CMB-S4/Euclid): $\Sigma m_\nu \approx 59\,\mathrm{meV}$, normal hierarchy.
  \item \textbf{No proton decay} (categorical): baryon number is the LLC on a topological winding number.
  \item \textbf{Strong CP}: $\theta_{\mathrm{QCD}} = 0$ exactly (DAG is acyclic; instantons cannot exist).
\end{enumerate}

%───────────────────────────────────────────────────────────────
\section{Open Problems}\label{sec:open}

The full derivation of $a_0^{\mathrm{RA}}$ (Eq.~\ref{eq:a0}) from the BDG diamond geometry, the perturbative Bianchi theorem on curved backgrounds (O03), and the Lean~4 formalization of the GR derivation (O10-GR) are open.

\section*{Acknowledgments}
Developed with AI collaboration (Claude, Anthropic; Gemini, Google).

\bibliographystyle{plain}
\begin{thebibliography}{99}
\bibitem{RA_P1} Sandeman, J.~F. (2026). Paper~I. Zenodo, doi:10.5281/zenodo.19362289.
\bibitem{RA_P5} Sandeman, J.~F. (2026). Paper~V. Zenodo, doi:10.5281/zenodo.19363077.
\bibitem{Weinberg1964} Weinberg, S. (1964). Photons and gravitons in S-matrix theory. \textit{Phys.\ Rev.}, 135, B1049.
\bibitem{Smolin1992} Smolin, L. (1992). Did the universe evolve? \textit{CQG}, 9, 173.
\bibitem{Milgrom1983} Milgrom, M. (1983). A modification of the Newtonian dynamics. \textit{ApJ}, 270, 365.
\bibitem{Durhuus1995} Durhuus, B., Jonsson, T., and Wheater, J. (1995). On the spectral dimension of random trees. \textit{J.\ Stat.\ Phys.}, 78, 827.
\bibitem{Planck2020} Planck Collaboration (2020). Planck 2018 results VI. \textit{A\&A}, 641, A6.
\bibitem{Jacobson1995} Jacobson, T. (1995). Thermodynamics of spacetime. \textit{Phys.\ Rev.\ Lett.}, 75, 1260.
\end{thebibliography}
\end{document}
